\section{\benchmark Design}
\label{sec:benchmark-design}

\benchmark is the final dataset produced by a configurable construction algorithm, which controls the method structure, the values that determine whether calls execute, and the presentation of the code. Figure~\ref{fig:pipeline} shows how this dataset is constructed and then used to evaluate model responses. The labels and executable graphs are part of the generated dataset, and their independent validation covers all its classes and questions, as described in Section~\ref{subsec:oracle-validation}.

\begin{figure*}[t]
\centering
\begin{tikzpicture}[rb diagram,
  pipeline card/.style={draw=rbRule, fill=white, line width=0.7pt,
    rounded corners=2pt, text width=2.75cm, minimum height=1.7cm,
    inner sep=6pt, align=center, font=\sffamily\footnotesize},
  pipeline prompt/.style={pipeline card, text width=3.15cm,
    minimum height=1.25cm, inner sep=5pt},
  pipeline dataset/.style={pipeline card, draw=modelblue!70,
    fill=modelblue!4, line width=0.9pt},
  pipeline link/.style={rb edge, rounded corners=2pt},
  pipeline reference/.style={pipeline link, dashed, draw=rbMuted}]

% The two rows read independently from left to right on the same grid.
\node[rb heading, anchor=west] at (0,1.25) {(a) Construction algorithm};
\node[pipeline card] (tree) at (1.6,0)
  {\textbf{Method structure}\\[3pt]Generated method names\\and call tree};
\node[pipeline card] (java) at (5.5,0)
  {\textbf{Java generation}\\[3pt]Values, calls and guards\\Fixed execution};
\node[pipeline card] (questions) at (9.4,0)
  {\textbf{Oracle + questions}\\[3pt]Executable calls\\Queries $(s,t)$\\Boolean labels};
\node[pipeline dataset, text width=3.35cm] (dataset) at (13.8,0)
  {\textbf{\benchmark}\\[3pt]Final generated dataset\\Classes, questions,\\labels and oracles};
\foreach \from/\to in {tree/java,java/questions,questions/dataset}
  \draw[pipeline link] (\from.east) -- (\to.west);
\node[rb note, anchor=west] at (0,-1.15)
  {Controlled settings: methods, guards, parameters, comments, variables and loops};

\node[rb heading, anchor=west] at (0,-1.7) {(b) Evaluation of the same dataset};
\node[pipeline dataset] (instance) at (1.6,-3.45)
  {\textbf{\benchmark} };
\node[pipeline prompt] (co) at (5.5,-2.65)
  {\textbf{CO: Class-only}\\[3pt]Entry point implicit\\Determinism implicit};
\node[pipeline prompt] (ee) at (5.5,-4.25)
  {\textbf{EE: Explicit execution}\\[3pt]Entry point: \texttt{Main}\\Determinism explicit};
\node[pipeline card] (model) at (9.4,-3.45)
  {\textbf{LLM inference}\\[3pt]One run per variant\\Boolean + trace};
\node[pipeline card, text width=3.35cm] (scoring) at (13.8,-3.45)
  {\textbf{Independent scoring}\\[3pt]Final-answer correctness\\Trace correctness};

% Local forks replace long routes around the construction stages.
\coordinate (input fork) at (3.55,-3.45);
\draw[pipeline link, -] (instance.east) -- (input fork);
\draw[pipeline link] (input fork) |- (co.west);
\draw[pipeline link] (input fork) |- (ee.west);
\coordinate (output join) at (7.55,-3.45);
\draw[pipeline link, -] (co.east) -| (output join);
\draw[pipeline link, -] (ee.east) -| (output join);
\draw[pipeline link] (output join) -- (model.west);
\draw[pipeline link] (model.east) -- (scoring.west);

% The oracle goes directly to scoring and never enters the prompt.
\draw[pipeline reference] (dataset.south) -- (scoring.north);
\node[rb annotation, anchor=west, align=left] at (14,-1.75)
  {Shared labels\\and oracle};
\end{tikzpicture}
\caption{Construction and evaluation of \benchmark. The algorithm produces the final dataset, whose classes and questions are presented with different execution information in CO and EE. The labels and executable-graph oracles are then used to score the returned answer and trace.}
\Description{Two aligned rows read from left to right. The construction algorithm generates a method tree, emits Java code, and constructs the oracle and questions, producing the final ReachabilityBench dataset. Evaluation takes the same class and question through either the Class-only or Explicit execution prompt, runs the model separately for each variant, and scores its Boolean answer and trace. A separate dashed arrow supplies the common labels and oracle directly to scoring.}
\label{fig:pipeline}
\end{figure*}


\subsection{Task Instances and Execution Context}

Each instance contains a Java class named \texttt{TheClass}, together with a source method $s$ and a target method $t$. The complete class is provided, including methods outside the region reachable from $s$. A root method with no parameters defines the entry point and initializes local values, which can then be passed to other methods during execution.

The execution from this root defines the context $e$, which supplies the argument values used to answer the question. The expected answer $y$ is \true{} when $t$ can execute from $s$ under these values, and \false{} otherwise. The source selected in the question is not necessarily the root, so its arguments may come from earlier calls. Each instance also retains the executable call graph $G_{\mathit{exec}}$ used as an oracle, together with the generation settings and question metadata.

Table~\ref{tab:generation-parameters} summarizes the configuration fields recorded with the instances and questions, distinguishing parameters that determine the class from those that control question sampling. Section~\ref{subsec:evaluation-set} gives the values used to construct the final dataset.

\begin{table}[!b]
\centering
\caption{Configuration fields and their role in instance generation and question sampling.}
\label{tab:generation-parameters}
\footnotesize
\begin{tabularx}{\columnwidth}{@{}>{\raggedright\arraybackslash}p{0.33\columnwidth}>{\raggedright\arraybackslash}X@{}}
\toprule
Parameter & Role \\
\midrule
\path{class_size} & Sets the number of methods in the class. \\
\path{depths} & Selects target positions in the generated exploration; negative questions use the corresponding negative distances. \\
\path{n_questions} & Sets the question quota for each polarity at each distance. \\
\path{n_params} & Sets the number of parameters for each non-root method; the root has no parameters. \\
\path{n_if} & Sets the number of condition blocks per method. Depending on the task family, these contain value-tested calls or control-flow filler. \\
\path{n_vars} & Adds artificial local variables, separate from variables introduced to pass call arguments. \\
\path{n_loops} & Adds loops per method. Each loop prints its counter for 1--5 iterations and contains no calls to benchmark methods. \\
\path{n_comment} & Adds Lorem Ipsum comment lines above each method body; they do not describe program behavior. \\
\path{useful_var_spawn_probability} & Controls the introduction of fresh local variables for call arguments when existing values could be reused. \\
\bottomrule
\end{tabularx}
\end{table}


\subsection{Constructing the Method Structure}

The algorithm first generates unique method names by combining a verb, a noun, and a suffix. This allows us to obtain identifiers resembling Java method names, although they do not specify what the generated bodies do. Each name then becomes a node in a single acyclic tree, whose size is controlled by \path{class_size}.

The tree is expanded using a queue, from which the algorithm takes a node $u$ and creates a left child $L(u)$ and a right child $R(u)$, as long as method names remain available. Additional branches can then be created below $L(u)$ by sampling its two possible children independently, each with probability $p=0.65$:
\[
    X_{LL}(u),X_{LR}(u)\sim\mathrm{Bernoulli}(p).
\]
The right child and any sampled grandchildren are added to the queue for further expansion, following the rule illustrated in Figure~\ref{fig:local-tree-rule}. This process stops when no method names remain. It produces reachable regions of different sizes, together with methods outside these regions that can be selected as negative targets.

\begin{figure}[t]
\centering
\begin{tikzpicture}[rb diagram]
\node[rb method] (u) at (0,0) {$u$};
\node[rb method] (l) at (-1.65,-1.05) {$L(u)$};
\node[rb method] (r) at (1.65,-1.05) {$R(u)$};
\node[rb method, minimum width=1.25cm] (ll) at (-2.5,-2.15) {$L(L(u))$};
\node[rb method, minimum width=1.25cm] (lr) at (-0.8,-2.15) {$R(L(u))$};
\draw[rb edge] (u) -- (l);
\draw[rb edge] (u) -- (r);
\draw[rb optional edge] (l) -- (ll);
\draw[rb optional edge] (l) -- (lr);
\node[rb annotation, anchor=east] at (-2,-1.55) {$X_{LL}(u)=1$};
\node[rb annotation, anchor=west] at (-1.15,-1.55) {$X_{LR}(u)=1$};
\draw[rb edge] (-2.8,-2.85) -- (-2.15,-2.85);
\node[rb note, anchor=west] at (-2.05,-2.85) {Fixed edge};
\draw[rb optional edge] (0,-2.85) -- (0.65,-2.85);
\node[rb note, anchor=west] at (0.75,-2.85) {Optional edge};
\node[rb note] at (-0.2,-3.3) {$X_{LL},X_{LR}\sim\mathrm{Bernoulli}(p)$, independently};
\end{tikzpicture}
\caption{Local tree-expansion rule. Solid edges are added while method names remain available, and dashed edges are sampled independently with probability $p$. The right child and sampled grandchildren are then queued for further expansion.}
\Description{A node u has two fixed children. The left child has two optional children controlled by independent Bernoulli variables.}
\label{fig:local-tree-rule}
\end{figure}


\subsection{Generating Methods and Propagating Values}

Once the tree is available, each node is translated into a public instance method returning \texttt{void}. The root has no parameters, while the other methods receive the number specified by \path{n_params}, with types selected from \texttt{int}, \texttt{long}, \texttt{boolean}, \texttt{float}, and \texttt{double}. When the parent method has parameters, the types used by its children are selected from those already present in the parent. This allows values to be passed through successive calls.

For each edge in the tree, the caller supplies arguments with the types required by the callee. It can reuse a compatible parameter or an argument variable already introduced, or create a local variable when needed. The parameter \path{useful_var_spawn_probability} controls whether a fresh variable is introduced when an existing value could be reused. Before generating the callee's body, the algorithm records the concrete value supplied to each parameter. The root therefore initializes values for its parameterized children, while later methods can also introduce new values. This means that a condition may depend on a value declared several methods earlier.

The parameter \path{n_vars} controls additional artificial local declarations, whose variables are not used as call arguments or to determine whether benchmark methods execute. Variables needed to supply arguments are generated separately and can affect guarded calls. As a consequence, setting \path{n_vars}=0 does not remove declarations needed to initialize execution, as illustrated in Figure~\ref{fig:background-example}.

\subsection{Conditions and Presentation Factors}

The parameter \mbox{\texttt{n\_if}} sets the number of condition blocks generated per method and distinguishes unguarded from condition-aware instances in the evaluation. When \mbox{\texttt{n\_if}$=0$}, calls corresponding to tree edges are unguarded. When \mbox{\texttt{n\_if}$>0$}, these calls can appear directly in the body or inside an \texttt{if} block, whose condition is constructed to evaluate to true with the recorded argument values.

A condition block that has not been used can instead contain calls to existing methods behind a false guard. The selected methods must accept arguments available in the caller and differ from both the caller and its direct executable callees. These calls remain visible in the code but are excluded from the executable graph, since their conditions prevent them from executing. They can connect separate branches or point back to an ancestor. The conditions use Boolean equality or inequality, or numeric comparisons with thresholds chosen to obtain the required result. When no suitable value is available, the algorithm uses literal \texttt{true} or \texttt{false} conditions. The resulting program therefore has a fixed execution, although its construction is randomized.

The parameters \path{n_comment} and \path{n_loops} introduce artificial elements that change the code's presentation without affecting reachability. The first controls the number of comment lines above each method body, which contain generated Lorem Ipsum text and do not describe its behavior. The second controls the number of \texttt{for} or \texttt{while} blocks in each body, which print a counter for one to five iterations and contain no calls to benchmark methods. Condition and loop blocks are shuffled inside methods, and method declarations are shuffled inside the class. Their position therefore does not reveal the call order.

The maintained generator also allows calls corresponding to tree edges to remain outside guards when condition blocks are present. In the evaluated archive, the condition-aware construction uses \mbox{\texttt{n\_if}}$\in\{0,1,3,5\}$. We therefore define the study groups using this count rather than the mode name stored in the archive.

\subsection{Constructing the Oracle and Selecting Questions}

While generating each Java body, the algorithm records the executable callees in the order in which their calls are emitted. It includes unguarded calls and calls behind true guards, while excluding calls behind false guards. These lists define $G_{\mathit{exec}}=(V,E_{\mathit{exec}})$ and retain the call order, since shuffling code blocks can change the reference exploration even when reachability remains the same.

The algorithm then explores this graph from candidate source methods to select question targets. For a positive question, the target is encountered during a source-first depth-first exploration. For a negative question, it is an existing method outside the complete region $R_e(s)$ reachable from the source. Both types of questions therefore refer to methods declared in the class, including when the target is unreachable.

The signed generation distance $d$ describes this reference exploration, rather than the shortest path between methods. For a positive question, it gives the target's position after the source, whose position is zero. For a negative question, $d=-(|R_e(s)|-1)$. A positive trace can therefore reach the target without visiting branches encountered earlier in the reference exploration, while a negative trace must cover the complete reachable region.

The values in \path{depths} determine which distances are sampled, while \path{n_questions} sets the quota for each polarity at each distance. Additional classes are generated when one class does not provide enough candidates to complete these quotas. The Java code is retained with the executable graph, method inventory, questions, and generation configuration. The maintained format stores the graph directly, while the evaluated archives retain executable trees from which the same relationships are recovered. These historical chain files do not impose a particular exploration order on the model.

\subsection{Prompting Variants of \benchmark}
\label{subsec:prompt-design}

The prompt contains a system message that explains executable reachability, specifies the JSON response, and presents the complete class containing its different methods. A separate user message names the source and target and asks whether the former calls the latter directly or indirectly. The model must return two fields: \texttt{final\_answer}, with the value \texttt{"True"} or \texttt{"False"}, and \texttt{trace}, listing the explored methods. Although the archived prompt asks for execution order, the evaluator accepts any exploration satisfying the rules in Section~\ref{subsec:scoring}.

We consider two variants of the final benchmark, \benchmark~(CO), for \emph{Class-only}, and \benchmark~(EE), for \emph{Explicit execution}. They use the same \texttt{TheClass} programs, source--target questions, Boolean labels, executable-graph oracles, and scoring rules. This allows us to compare two presentations of the same dataset, which differ in the execution information made explicit in the prompt.

In \benchmark~(CO), the code block contains only \texttt{TheClass}, without explicitly identifying the entry point or stating that all relevant values and execution outcomes are determined by the code. When the source has parameters, the model must infer the intended execution from the root and recover the values supplied to this source. A parameter may therefore be treated as an unspecified input, even though its value is fixed by the generated execution. This introduces an additional interpretation step, particularly when the arguments determine whether guarded calls execute.

In \benchmark~(EE), the code block also contains a separate \texttt{Main} launcher, which creates an instance of \texttt{TheClass} and calls its root without arguments. This launcher makes the entry point visible but is not included in the oracle graph. The prompt explains that the code describes a concrete execution, that all relevant values and outcomes can be determined statically, and that no external inputs should be assumed. The model must therefore follow values from the entry point and use them to evaluate conditions, branches, and loops.

EE therefore makes explicit the execution context already used by the common oracle, while CO requires the model to infer it. The structural complexity of the classes remains the same, and the returned trace starts at the queried source in both variants. \researchq{3} examines how performance and token use differ between these two presentations of the same dataset.
